\chapter{不决定神经元类型的物质}
\label{sec:othermediators}

在第~\ref{sec:neurontypes}~章中，我们按主要神经递质对神经元分类，得到了十种类型（表~\ref{tab:types}）。现在该加上一些复杂情况了。神经元释放的不只是主要递质，而且在大脑中，除了主要递质，还有别的物质在起作用。它们改变网络传递和存储信号的方式。

除了寻址总线和广播总线（第~\ref{sec:buses}~节），还有第三种：\emph{反向通道}。有几种物质不是由发送者、而是由\emph{接收者}释放的，它们穿过突触间隙向回走，到达发送者的输出端，改变它下一次释放多少递质。在通信网络中，这类似于接收确认，发送者用它来调整自己的发送。神经元在改变连接权重时用的正是这个通道：发送者的输出端通过它得知接收者的活动，也就是第~\ref{sec:dofa}~章学习规则中的第二个因子。

已知神经递质的完整清单很长：一百多种物质，其中大多数是短的蛋白质链，即神经肽。但它们都可以归入少数几类，同一类中的物质行为相似。表~\ref{tab:othermediators} 收集了从不充当主要递质的物质，并回答工程师会问的同样三个问题：信号走哪条总线，作用快不快，通常是什么符号。它们不决定神经元的类型：它们要么与主要递质一起释放，要么不是由神经元释放的，要么是反向传递的。

\begin{table}[t]
\centering
\footnotesize
\setlength{\tabcolsep}{4pt}
\renewcommand{\arraystretch}{1.12}
\begin{tabular}{>{\raggedright}p{2.25cm}>{\raggedright}p{2.8cm}>{\raggedright}p{1.8cm}>{\raggedright}p{1.5cm}>{\centering}p{0.8cm}>{\raggedright\arraybackslash}p{4.3cm}}
\toprule
类别 & 物质 & 总线 & 速度 & 符号 & 在计算中的作用 \\
\midrule
氨基酸 & D-丝氨酸 & 局部 & 慢 & 与 & 谷氨酸受体的第二把钥匙：“与”条件 \\
\midrule
单胺 & 痕量胺 & 广播 & 慢 & $\pm$ & 单胺通道的微调 \\
\midrule
\multirow{2}{2.25cm}{嘌呤}
 & ATP & 寻址 & 快和慢 & $+$ & 共释放递质；神经元与胶质细胞之间的信号 \\
 & 腺苷 & 容积 & 慢 & $-$ & 随活动积累：负荷计数器~\cite{Dunwiddie2001} \\
\midrule
\multirow{8}{2.25cm}{神经肽（一百多种）}
 & 脑啡肽、内啡肽、强啡肽 & 容积 & 慢 & $-$ & 抑制传递 \\
 & P 物质 & 容积 & 慢 & $+$ & 缓慢增强 \\
 & 神经肽 Y & 容积 & 慢 & $-$ & 缓慢抑制 \\
 & 生长抑素 & 容积 & 慢 & $-$ & 缓慢抑制，与 GABA 共释放 \\
 & 血管活性肠肽（VIP） & 容积 & 慢 & $+$ & 去抑制：抑制那些抑制性神经元 \\
 & 胆囊收缩素 & 容积 & 慢 & $\pm$ & 在部分抑制性神经元中与 GABA 共释放 \\
 & 食欲素 & 广播 & 慢 & $+$ & 觉醒状态的稳定器 \\
 & 催产素、加压素 & 广播 & 慢 & $\pm$ & 行为的全局设置 \\
\midrule
\multirow{2}{2.25cm}{气体}
 & 一氧化氮 NO & 反向、容积 & 慢 & $\pm$ & 能穿过细胞膜；权重改变时的反向信号~\cite{Garthwaite2008} \\
 & CO、H$_2$S & 容积 & 慢 & $\pm$ & 作用研究很少 \\
\midrule
脂质 & 内源性大麻素（花生四烯酸乙醇胺、2-AG） & 反向 & 慢 & $-$ & 接收者减少发送者的释放：按权重的反馈~\cite{WilsonNicoll2001} \\
\bottomrule
\end{tabular}
\caption{不决定神经元类型的物质。\emph{总线}、\emph{速度}和\emph{符号}的含义同表~\ref{tab:types}；此外，容积总线指递质在释放点周围的细胞间隙中扩散，反向总线指从接收者回到发送者，局部总线指在释放点附近起作用。“与”是许可信号：它本身不引起电流，但没有它，另一个输入就不起作用，就像“与”逻辑门的第二个输入。神经肽几乎总是与主要递质一起释放，并且主要在脉冲频率较高时释放~\cite{vandenPol2012}。}
\label{tab:othermediators}
\end{table}
